%---------------------------------------------------------------------
% Brief edition: condensed formula sheet and the eight mistakes.
%---------------------------------------------------------------------

\section*{فارمولا شیٹ}
\markboth{فارمولا شیٹ}{}

\subsection*{احتمال اور بے ترتیب متغیر (یونٹ \textenglish{1--2})}

\begin{center}
\begin{tabular}{@{}p{3.6cm} p{6.0cm} p{5.4cm}@{}}
\toprule
\textbf{اصول} & \textbf{فارمولا} & \textbf{کب} \\
\midrule
متمم        & \(\prob{A'}=1-\prob{A}\) & ''کم از کم ایک''، ''نہیں'' \\
جمع (عام)   & \(\prob{A\cup B}=\prob{A}+\prob{B}-\prob{A\cap B}\) & ''یا''، اشتراک موجود \\
جمع (خارج)  & \(\prob{A\cup B}=\prob{A}+\prob{B}\) & ''یا''، \(A\cap B=\varnothing\) \\
ضرب (آزاد)  & \(\prob{A\cap B}=\prob{A}\prob{B}\) & ''اور''، واپسی کے ساتھ \\
ضرب (وابستہ)& \(\prob{A\cap B}=\prob{A}\prob{B\mid A}\) & ''اور''، بغیر واپسی \\
مشروط       & \(\prob{A\mid B}=\dfrac{\prob{A\cap B}}{\prob{B}}\) & ''بشرطیکہ'' \\
آزادی کی جانچ & \(\prob{A\cap B}=\prob{A}\prob{B}\) & آزادی ثابت یا رد کرنے کو \\
دو ذی حدی   & \(P(X=x)=\dbinom{n}{x}p^{x}(1-p)^{n-x},\ \ E(X)=np\) & ثابت \(p\)، واپسی کے ساتھ \\
تقسیم کی شرط & \(\sum_x P(X=x)=1,\ \ E(X)=\sum_x x\,P(X=x)\) & ہمیشہ --- پڑتال کے طور پر \\
\bottomrule
\end{tabular}
\end{center}

\subsection*{مساوات اور خطوط (یونٹ \textenglish{3--4})}

\begin{center}
\begin{tabular}{@{}p{4.0cm} p{11.0cm}@{}}
\toprule
دو درجی فارمولا & \(x=\dfrac{-b\pm\sqrt{b^{2}-4ac}}{2a}\)؛ ممیز \(>0\) دو جڑیں، \(=0\) ایک، \(<0\) کوئی حقیقی نہیں \\
عدم مساوات & منفی عدد سے ضرب یا تقسیم پر علامت \emph{الٹ} جاتی ہے \\
مطلق قدر & \(|u|=|v| \iff u=\pm v\)؛ \ \(|u|\le a \iff -a\le u\le a\) (ایک وقفہ)؛ \ \(|u|\ge a\) (دو وقفے) \\
فاصلہ، وسطی نقطہ & \(d=\sqrt{(x_2-x_1)^{2}+(y_2-y_1)^{2}}\)، \ \(M=\left(\tfrac{x_1+x_2}{2},\tfrac{y_1+y_2}{2}\right)\) \\
میلان اور خط & \(m=\dfrac{y_2-y_1}{x_2-x_1}\)، \ \(y-y_1=m(x-x_1)\)، \ \(y=mx+c\) \\
متوازی / عمود & \(m_1=m_2\)؛ \ \(m_1m_2=-1\) --- عمودی خطوط پر لاگو نہیں \\
خطی لاگت & \(C(x)=(\text{فی اکائی متغیر})x+(\text{ثابت لاگت})\) \\
توازن & طلب \(=\) رسد رکھ کر \(q\) نکالیے، پھر \(p\) میں تعویض \\
\bottomrule
\end{tabular}
\end{center}

\subsection*{میٹرکس اور محدد (یونٹ \textenglish{5--6})}

\begin{center}
\begin{tabular}{@{}p{4.0cm} p{11.0cm}@{}}
\toprule
ابعاد کا میل & \(A_{m\times n}B_{n\times p}\) ممکن، نتیجہ \(m\times p\) \\
ترتیب الٹتی ہے & \(\boxed{(AB)^{t}=B^{t}A^{t}}\)، \ \(\boxed{(AB)^{-1}=B^{-1}A^{-1}}\) \\
\(2\times2\) & \(\dt{a&b\\c&d}=ad-bc\)، \ معکوس \(=\dfrac{1}{ad-bc}\mat{d&-b\\-c&a}\) \\
ہم عامل & \(C_{ij}=(-1)^{i+j}M_{ij}\)؛ علامتیں \(\mat{+&-&+\\-&+&-\\+&-&+}\) \\
محدد & \(\det A=\sum_j a_{ij}C_{ij}\) کسی بھی سطر یا ستون سے \\
ملحقہ اور معکوس & \(\operatorname{adj}A=\left[C_{ij}\right]^{t}\)، \ \(A^{-1}=\dfrac{\operatorname{adj}A}{\det A}\)، شرط \(\det A\neq0\) \\
کریمر کا اصول & \(x_k=\dfrac{D_k}{D}\)؛ \(D_k\) میں \(k\)واں ستون \(b\) سے بدلیں \\
\(\det=0\) کا مطلب & منفرد حل نہیں: یا کوئی حل نہیں، یا لامتناہی \\
مختصر راستے & زیادہ صفر والی سطر سے پھیلائیں؛ مشترک عامل باہر؛ دو یکساں سطریں \(\Rightarrow \det=0\) \\
\bottomrule
\end{tabular}
\end{center}

\subsection*{کیلکولس (یونٹ \textenglish{7--9})}

\begin{center}
\begin{tabular}{@{}p{4.0cm} p{11.0cm}@{}}
\toprule
قوت، حاصل ضرب & \(\dd{}{x}x^{n}=nx^{n-1}\)، \ \((uv)'=u'v+uv'\) \\
تقسیم & \(\left(\dfrac uv\right)'=\dfrac{u'v-uv'}{v^{2}}\) --- \emph{اوپر} والے کا مشتق پہلے \\
زنجیری & \(\bigl(f(g(x))\bigr)'=f'(g(x))\,g'(x)\) --- باہر پہلے، پھر اندر سے ضرب \\
تفاعلی، لاگ & \(\dd{}{x}e^{u}=e^{u}u'\)، \ \(\dd{}{x}\ln u=\dfrac{u'}{u}\) \\
حدیں & \(\infty\) پر: نسب نما کی بڑی قوت سے تقسیم۔ \(\tfrac00\) پر: تحلیل و منسوخی \\
حد کا وجود & \(\lim_{x\to a^{-}}f=\lim_{x\to a^{+}}f\) \\
جزوی مشتق & ایک متغیر میں تفرق، دوسرا ثابت؛ \ \(f_{xy}=f_{yx}\) (مفت پڑتال) \\
جانچ (ایک متغیر) & \(f'=0\) پر: \(f''<0\) اعظم، \(f''>0\) اصغر، \(f''=0\) بے نتیجہ \\
جانچ (دو متغیر) & \(D=f_{xx}f_{yy}-(f_{xy})^{2}\)؛ \(D>0\) و \(f_{xx}<0\) اعظم، \(f_{xx}>0\) اصغر؛ \(D<0\) زینی \\
انعطافی نقطہ & \(f''=0\) \emph{اور} \(f''\) کی علامت بدلے \\
کاروباری فعل & \(R=x\,p(x)\)، \ \(P=R-C\) (اعظم جہاں \(MR=MC\))، \ \(\bar{C}=\tfrac{C}{x}\) (اصغر پر \(MC=\bar{C}\)) \\
\bottomrule
\end{tabular}
\end{center}

\clearpage

%---------------------------------------------------------------------
\section*{آٹھ عام غلطیاں}
\markboth{آٹھ عام غلطیاں}{}

\begin{enumerate}\setlength{\itemsep}{5pt}
  \item \textbf{اشتراک والے واقعات پر \(\prob{A}+\prob{B}\) لگانا۔} سادہ جمع کا
        اصول صرف باہم خارج واقعات پر چلتا ہے؛ ورنہ اشتراک گھٹائیے۔
  \item \textbf{''صرف ایک'' میں گنتی بھول جانا۔} \(\binom{n}{1}\) یہ گنتا ہے کہ
        وہ ایک \emph{کون} ہے۔ ''سب'' اور ''کوئی نہیں'' میں یہ ضرب نہیں لگتی۔
  \item \textbf{''بغیر واپسی'' پر دو ذی حدی لگا دینا۔} یہ الفاظ
        \textenglish{hypergeometric} کی طرف اشارہ ہیں اور جواب بدل دیتے ہیں۔
  \item \textbf{منفی عدد سے ضرب یا تقسیم پر عدم مساوات کی علامت نہ الٹنا۔} نیز
        \(|u|\le a\) ایک وقفہ دیتا ہے اور \(|u|\ge a\) دو۔
  \item \textbf{\((AB)^{t}=A^{t}B^{t}\) یا \((AB)^{-1}=A^{-1}B^{-1}\) لکھنا۔}
        دونوں میں ترتیب الٹ جاتی ہے۔
  \item \textbf{ملحقہ بناتے وقت ہم عاملوں کا منقلب بھول جانا۔} پڑتال کسی
        \emph{غیر قطری} خانے پر کیجیے، صرف قطری پر نہیں۔
  \item \textbf{کریمر لگانے سے پہلے \(D\) نہ نکالنا۔} \(D=0\) پر اصول لاگو نہیں
        ہوتا۔ نیز غائب متغیرات کے صفر عامل پہلے لکھ لیجیے۔
  \item \textbf{\(f'=0\) پر رک جانا۔} دوسرے مشتق کی جانچ ہی اعظم، اصغر اور زینی
        نقطے میں فرق کرتی ہے، اور اس کے نمبر الگ ملتے ہیں۔
\end{enumerate}

\vfill
\begin{center}
\begin{tcolorbox}[enhanced, colback=bmtint, colframe=bmblue, boxrule=0.5pt,
                  arc=3pt, width=0.94\linewidth,
                  left=11pt, right=11pt, top=9pt, bottom=9pt]
\small
\textbf{\color{bmblue}اگر وقت کم ہو۔}\ \
یہ پانچ مہارتیں حل ڈھانپ کر دہرائیے:
\textenglish{(1)} \(3\times3\) نظام پر کریمر کا اصول؛
\textenglish{(2)} ہم عاملوں سے \(3\times3\) کا معکوس؛
\textenglish{(3)} حاصل ضرب، تقسیم اور زنجیری اصول؛
\textenglish{(4)} منافع یا اوسط لاگت کی بہتر بندی، دوسرے مشتق کی جانچ سمیت؛
\textenglish{(5)} دو نقطوں یا عمودیت کی شرط سے خط کی مساوات۔
پانچ مہارتیں جو بغیر دیکھے درست آ جائیں، نو یونٹ سرسری پڑھ لینے سے زیادہ قیمتی
ہیں۔

\medskip
\textbf{\color{bmblue}مزید تفصیل۔}\ \
\textenglish{27} سوالات کا اردو ایڈیشن
\textenglish{AIOU\_1429\_Urdu\_Guess\_Paper.pdf}؛ مکمل انگریزی حل
\textenglish{AIOU\_1429\_Solved\_QuestionBank.pdf}۔
\ \textbf{ماخذ:} \emph{کاروباری ریاضی}، کورس کوڈ \textenglish{5405/1429}،
یونٹ \textenglish{1--9}، اے آئی او یو اسلام آباد، اشاعت \textenglish{2023}۔
\end{tcolorbox}
\end{center}
